\chapter{Triángulos y cuadriláteros}\label{ch:g6:shapes}

Los triángulos y los \omterm{def:g4:geometry:polygon}{cuadriláteros} son las figuras básicas de la geometría. El objetivo
de este capítulo es conocer sus familias especiales --- \omterm{def:g6:shapes:triangles}{isósceles},
\omterm{def:g6:shapes:triangles}{equiláteros}, \omterm{def:g6:shapes:triangles}{triángulos rectángulos}; rectángulos, \omterm{def:g6:shapes:quadrilaterals}{rombos}, cuadrados --- por
sus \emph{propiedades características} y saber construirlos con precisión con
regla y compás.

\section{Los triángulos}

\begin{definition}[Triángulos especiales]\label{def:g6:shapes:triangles}
Un triángulo $ABC$ tiene tres \omterm{def:g5:solids:def}{vértices} y tres lados. Es:
\begin{itemize}
  \item \emph{isósceles}\index{triángulo isósceles} en $A$ cuando
        $AB = AC$ (dos lados iguales);
  \item \emph{equilátero}\index{triángulo equilátero} cuando
        $AB = BC = CA$ (tres lados iguales);
  \item \emph{rectángulo}\index{triángulo rectángulo} en $A$ cuando el ángulo
        $\widehat{BAC}$ es un \omterm{def:g3:shapes:rightangle}{ángulo recto}. El lado $[BC]$, opuesto al
        \omterm{def:g3:shapes:rightangle}{ángulo recto}, se llama \emph{hipotenusa}.
\end{itemize}
\end{definition}

\begin{omfigure}
\begin{tikzpicture}[scale=0.85]
  % isosceles
  \draw[thick] (0,0) -- (2.4,0) -- (1.2,2.1) -- cycle;
  \draw (0.55,0.96) -- (0.72,0.86);
  \draw (1.85,0.96) -- (1.68,0.86);
  \node[below, font=\small] at (1.2,-0.15) {isósceles};
  \node[above, font=\small] at (1.2,2.1) {$A$};
  \node[below left, font=\small] at (0,0) {$B$};
  \node[below right, font=\small] at (2.4,0) {$C$};
  % equilatero
  \begin{scope}[xshift=4.6cm]
  \draw[thick] (0,0) -- (2.4,0) -- (1.2,2.08) -- cycle;
  \draw (1.13,0.05) -- (1.27,-0.05);
  \draw (0.55,0.99) -- (0.72,0.89);
  \draw (1.85,0.99) -- (1.68,0.89);
  \node[below, font=\small] at (1.2,-0.25) {equilátero};
  \end{scope}
  % rectangulo
  \begin{scope}[xshift=9.2cm]
  \draw[thick] (0,0) -- (2.6,0) -- (0,1.8) -- cycle;
  \draw[thin] (0.3,0) -- (0.3,0.3) -- (0,0.3);
  \node[below, font=\small] at (1.3,-0.15) {rectángulo};
  \node[omThm, font=\small, rotate=-34] at (1.5,1.05) {hipotenusa};
  \end{scope}
\end{tikzpicture}

{\small Los tres triángulos especiales. Los trazos pequeños indican lados
iguales; el cuadradito marca el \omterm{def:g3:shapes:rightangle}{ángulo recto}.}
\end{omfigure}

\begin{method}[Construir un triángulo a partir de sus tres lados]\label{met:g6:shapes:construct}
Para construir $ABC$ con $AB = 5$ cm, $AC = 4$ cm y $BC = 3$ cm:
\begin{enumerate}
  \item traza el \omterm{def:g6:lines:objects}{segmento} $[AB]$ de $5$ cm con la regla;
  \item traza un arco de la circunferencia de centro $A$ y \omterm{def:g3:shapes:shapes}{radio} $4$ cm
        (todos los puntos a $4$ cm de $A$);
  \item traza un arco de la circunferencia de centro $B$ y \omterm{def:g3:shapes:shapes}{radio} $3$ cm;
  \item el punto $C$ es donde se cortan los dos arcos; únelo con $A$ y
        $B$.
\end{enumerate}
\end{method}

\begin{omfigure}
\begin{tikzpicture}[scale=1.0]
  \coordinate (A) at (0,0);
  \coordinate (B) at (5,0);
  \coordinate (C) at (3.2,2.4);
  \draw[thick] (A) -- (B);
  \draw[thick, omDef] (A) -- (C);
  \draw[thick, omThm] (B) -- (C);
  \draw[omDef, thin] (C) arc[start angle=36.9, delta angle=22, radius=4];
  \draw[omDef, thin] (C) arc[start angle=36.9, delta angle=-22, radius=4];
  \draw[omThm, thin] (C) arc[start angle=126.9, delta angle=22, radius=3];
  \draw[omThm, thin] (C) arc[start angle=126.9, delta angle=-22, radius=3];
  \fill (A) circle (1.5pt) node[below left, font=\small] {$A$};
  \fill (B) circle (1.5pt) node[below right, font=\small] {$B$};
  \fill (C) circle (1.5pt) node[above, font=\small] {$C$};
  \node[omDef, font=\small] at (1.2,1.5) {$4$ cm};
  \node[omThm, font=\small] at (4.6,1.4) {$3$ cm};
  \node[below, font=\small] at (2.5,0) {$5$ cm};
\end{tikzpicture}

{\small Construcción de un triángulo con el compás: el \omterm{def:g5:solids:def}{vértice} $C$ es el
punto de corte de los dos arcos. (En el \cref{ch:g7:triangles} verás
cuándo es posible una construcción así.)}
\end{omfigure}

\section{Los cuadriláteros}

\begin{definition}[Cuadriláteros especiales]\label{def:g6:shapes:quadrilaterals}
Un \omterm{def:g4:geometry:polygon}{cuadrilátero} tiene cuatro \omterm{def:g5:solids:def}{vértices} y cuatro lados. Es un:
\begin{itemize}
  \item \emph{rectángulo}\index{rectángulo} cuando sus cuatro ángulos son
        rectos;
  \item \emph{rombo}\index{rombo} cuando sus cuatro lados tienen la misma
        longitud;
  \item \emph{cuadrado}\index{cuadrado} cuando es las dos cosas: cuatro \omterm{def:g3:shapes:rightangle}{ángulos rectos}
        \emph{y} cuatro lados iguales.
\end{itemize}
\end{definition}

\begin{proposition}[Propiedades de las diagonales]\label{prop:g6:shapes:diagonals}
En estos \omterm{def:g4:geometry:polygon}{cuadriláteros}, las dos diagonales se cortan en su
punto medio común y, además:
\begin{itemize}
  \item en un rectángulo, las diagonales tienen la \emph{misma longitud};
  \item en un \omterm{def:g6:shapes:quadrilaterals}{rombo}, las diagonales son \emph{\omterm{def:g6:lines:perp}{perpendiculares}};
  \item en un cuadrado, las dos cosas a la vez.
\end{itemize}
\end{proposition}
\admitted

\begin{omfigure}
\begin{tikzpicture}[scale=0.8]
  % rectangulo, marcas de angulo recto hacia dentro en cada esquina
  \draw[thick] (0,0) rectangle (3,1.9);
  \draw[omThm] (0,0) -- (3,1.9);
  \draw[omThm] (0,1.9) -- (3,0);
  \foreach \x/\y/\dx/\dy in {0/0/1/1, 3/0/-1/1, 0/1.9/1/-1, 3/1.9/-1/-1}
    \draw[thin] (\x+0.22*\dx,\y) -- (\x+0.22*\dx,\y+0.22*\dy)
      -- (\x,\y+0.22*\dy);
  \node[below, font=\small] at (1.5,-0.2) {rectángulo};
  % rombo
  \begin{scope}[xshift=5.1cm, yshift=0.95cm]
  \draw[thick] (0,0) -- (1.6,1.0) -- (3.2,0) -- (1.6,-1.0) -- cycle;
  \draw[omThm] (0,0) -- (3.2,0);
  \draw[omThm] (1.6,1.0) -- (1.6,-1.0);
  \draw[thin] (1.6,0) ++(0.18,0) -- ++(0,0.18) -- ++(-0.18,0);
  \node[below, font=\small] at (1.6,-1.2) {rombo};
  \end{scope}
  % cuadrado, tambien con sus marcas de angulo recto
  \begin{scope}[xshift=9.8cm]
  \draw[thick] (0,0) rectangle (1.9,1.9);
  \draw[omThm] (0,0) -- (1.9,1.9);
  \draw[omThm] (0,1.9) -- (1.9,0);
  \foreach \x/\y/\dx/\dy in {0/0/1/1, 1.9/0/-1/1, 0/1.9/1/-1,
                             1.9/1.9/-1/-1}
    \draw[thin] (\x+0.22*\dx,\y) -- (\x+0.22*\dx,\y+0.22*\dy)
      -- (\x,\y+0.22*\dy);
  \node[below, font=\small] at (0.95,-0.2) {cuadrado};
  \end{scope}
\end{tikzpicture}

{\small Los tres \omterm{def:g4:geometry:polygon}{cuadriláteros} especiales con sus diagonales (en rojo).
El cuadrado hereda las propiedades del rectángulo y las del
\omterm{def:g6:shapes:quadrilaterals}{rombo}.}
\end{omfigure}

\begin{remark}
Un cuadrado es automáticamente un rectángulo (tiene cuatro \omterm{def:g3:shapes:rightangle}{ángulos rectos}) y
automáticamente un \omterm{def:g6:shapes:quadrilaterals}{rombo} (cuatro lados iguales). ¡Así que «rectángulo» no
significa «que no es cuadrado»! En matemáticas, la figura más especial pertenece
a la familia más general.
\end{remark}

\begin{example}[Reconocer una figura por sus propiedades]\label{ex:g6:shapes:recognize}
Un \omterm{def:g4:geometry:polygon}{cuadrilátero} tiene cuatro lados de $4$ cm y un \omterm{def:g3:shapes:rightangle}{ángulo recto}.
¿Qué es? Los cuatro lados iguales lo convierten en un \omterm{def:g6:shapes:quadrilaterals}{rombo}. En un \omterm{def:g6:shapes:quadrilaterals}{rombo}, los ángulos
opuestos son iguales y los cuatro ángulos suman $360^\circ$; un \omterm{def:g3:shapes:rightangle}{ángulo
recto} obliga entonces a que los cuatro lo sean (se demostrará con los
paralelogramos en el \cref{ch:g7:central}). Así que la figura es un cuadrado.
\end{example}

\begin{method}[Programas de construcción]\label{met:g6:shapes:program}
Un \emph{programa de construcción} es una lista de instrucciones que cualquiera puede
seguir para reproducir una figura exactamente. Escribe una instrucción por línea,
nombra cada punto en cuanto se cree y usa solo: «traza el \omterm{def:g6:lines:objects}{segmento}
/ la recta / la circunferencia\dots», «coloca el punto donde se cortan \dots». Prueba
tu programa siguiéndolo al pie de la letra.
\end{method}

\begin{example}\label{ex:g6:shapes:program}
Programa para un cuadrado $ABCD$ de lado $4$ cm:
\begin{enumerate}
  \item traza un \omterm{def:g6:lines:objects}{segmento} $[AB]$ con $AB = 4$ cm;
  \item traza la perpendicular a $(AB)$ en $A$; sobre ella, coloca $D$ a
        $4$ cm de $A$ (elige un lado de la recta y quédate en él);
  \item traza la perpendicular a $(AB)$ en $B$; sobre ella, coloca $C$ a
        $4$ cm de $B$, del mismo lado que $D$;
  \item une $C$ y $D$.
\end{enumerate}
Comprueba con la regla: $DC$ debería medir también $4$ cm.
\end{example}

\section{Ejercicios}

\begin{exercise}[$\star$]\label{exo:g6:shapes:1}
Construye un triángulo $ABC$ con $AB = 6$ cm, $AC = 5$ cm y
$BC = 4$ cm, siguiendo el \cref{met:g6:shapes:construct}.
\end{exercise}

\begin{exercise}[$\star$]\label{exo:g6:shapes:2}
Construye un \omterm{def:g6:shapes:triangles}{triángulo isósceles} $DEF$ con $DE = DF = 5$ cm y
$EF = 3$ cm. ¿Qué dos ángulos de este triángulo parecen iguales? Mídelos
para comprobarlo.
\end{exercise}

\begin{exercise}[$\star$]\label{exo:g6:shapes:3}
Construye un \omterm{def:g6:shapes:triangles}{triángulo equilátero} de lado $4$ cm usando solo regla y
compás. Mide uno de sus ángulos.
\end{exercise}

\begin{exercise}[$\star$]\label{exo:g6:shapes:4}
Clasifica cada triángulo a partir de sus datos (\omterm{def:g6:shapes:triangles}{isósceles}, \omterm{def:g6:shapes:triangles}{equilátero} o
rectángulo --- puede valer más de una respuesta):
(a) lados $7$, $7$, $7$; \
(b) lados $5$, $5$, $8$; \
(c) ángulos $90^\circ$, $45^\circ$, $45^\circ$; \
(d) ángulos $60^\circ$, $60^\circ$, $60^\circ$.
\end{exercise}

\begin{exercise}[$\star$]\label{exo:g6:shapes:5}
En un \omterm{def:g6:shapes:triangles}{triángulo rectángulo}, ¿qué lado es la hipotenusa? Dibuja un triángulo
$MNP$ rectángulo en $N$ y nombra su hipotenusa. ¿Qué \omterm{def:g5:solids:def}{vértice} no
toca nunca la hipotenusa?
\end{exercise}

\begin{exercise}[$\star$]\label{exo:g6:shapes:6}
¿Verdadero o falso? Justifica cada respuesta en una frase.
\begin{enumerate}
  \item todo cuadrado es un rectángulo;
  \item todo rectángulo es un cuadrado;
  \item todo cuadrado es un \omterm{def:g6:shapes:quadrilaterals}{rombo};
  \item un \omterm{def:g6:shapes:quadrilaterals}{rombo} con un \omterm{def:g3:shapes:rightangle}{ángulo recto} es un cuadrado.
\end{enumerate}
\end{exercise}

\begin{exercise}[$\star$]\label{exo:g6:shapes:7}
Construye un rectángulo $ABCD$ con $AB = 5$ cm y $BC = 3$ cm. Mide
sus dos diagonales: ¿qué observas?
\end{exercise}

\begin{exercise}[$\star$]\label{exo:g6:shapes:8}
Construye un \omterm{def:g6:shapes:quadrilaterals}{rombo} con diagonales de $6$ cm y $4$ cm. (Usa la
\cref{prop:g6:shapes:diagonals}: traza primero las dos diagonales ---
\omterm{def:g6:lines:perp}{perpendiculares} y cortándose en sus puntos medios --- y une después los cuatro
extremos.)
\end{exercise}

\begin{exercise}[$\star\star$]\label{exo:g6:shapes:9}
Escribe un programa de construcción para un \omterm{def:g6:shapes:triangles}{triángulo isósceles} de base
$4$ cm y lados iguales de $6$ cm, al estilo del
\cref{ex:g6:shapes:program}. Intercámbialo con un compañero: ¿puede seguirlo
sin ver tu figura?
\end{exercise}

\begin{exercise}[$\star\star$]\label{exo:g6:shapes:10}
Traza una circunferencia de centro $O$ y uno de sus \omterm{def:g4:geometry:circle}{diámetros} $[AB]$. Sitúa un
punto cualquiera $M$ sobre la circunferencia y dibuja el triángulo $AMB$. Mide el ángulo
$\widehat{AMB}$. Mueve $M$ y vuelve a medir. ¿Qué conjeturas?
(Este hecho sorprendente se demuestra en el \cref{ch:g8:cosine}.)
\end{exercise}

\begin{exercise}[$\star\star$]\label{exo:g6:shapes:11}
Un \omterm{def:g4:geometry:polygon}{cuadrilátero} tiene las diagonales de la misma longitud y se cortan en
su punto medio común. Averigua qué figura tiene que ser, usando la
\cref{prop:g6:shapes:diagonals}: ¿rectángulo, \omterm{def:g6:shapes:quadrilaterals}{rombo} o cuadrado?
(Cuidado: solo se da una propiedad.)
\end{exercise}

\begin{exercise}[$\star\star\star$]\label{exo:g6:shapes:12}
En una trama de $4 \times 4$ puntos, ¿cuántos cuadrados distintos puedes dibujar
con los \omterm{def:g5:solids:def}{vértices} en los puntos de la trama --- contando también los cuadrados «inclinados»?
Dibuja al menos un cuadrado inclinado y explica por qué sus lados son iguales.
\end{exercise}

\section{Problema: Lo que saben las diagonales}

\begin{problem}[{Problema de fin de semana --- reconstruir cada cuadrilátero
especial a partir únicamente de sus diagonales y contar las diagonales
de un polígono cualquiera}]
\label{pb:g6:shapes:1}
La \Cref{prop:g6:shapes:diagonals} lee una figura y te habla de
sus diagonales. Este problema juega \emph{al revés}: dados
solo dos \omterm{def:g6:lines:objects}{segmentos} que se cortan --- las futuras diagonales ---, ¿qué figura
aparece al unir sus cuatro extremos? Construirás un
«diccionario de diagonales» completo, entenderás \emph{por qué} funciona la
construcción del \omterm{def:g6:shapes:quadrilaterals}{rombo} del \cref{exo:g6:shapes:8}, conocerás una figura que
todavía no tiene nombre y terminarás contando las diagonales de \omterm{def:g4:geometry:polygon}{polígonos}
con muchos más de cuatro lados.

\medskip
\noindent\textbf{Parte I --- El diccionario de las diagonales.}
Para cada pregunta, traza los dos \omterm{def:g6:lines:objects}{segmentos} como se describe, une sus
cuatro extremos en orden e identifica el \omterm{def:g4:geometry:polygon}{cuadrilátero} que obtienes.
\begin{enumerate}
  \item dos \omterm{def:g6:lines:objects}{segmentos} de $6$ cm cada uno, \omterm{def:g6:lines:perp}{perpendiculares}, que se cortan en
        su punto medio común. Mide los cuatro lados y los
        cuatro ángulos de la figura resultante: ¿qué es?
  \item dos \omterm{def:g6:lines:objects}{segmentos} de $8$ cm y $5$ cm, \omterm{def:g6:lines:perp}{perpendiculares},
        que se cortan en su punto medio común. ¿Qué figura aparece
        (\cref{exo:g6:shapes:8})?
  \item dos \omterm{def:g6:lines:objects}{segmentos} de $7$ cm cada uno --- iguales pero \emph{no}
        \omterm{def:g6:lines:perp}{perpendiculares} --- que se cortan en su punto medio común.
        Mide los lados y los ángulos: ¿qué conjeturas que es la
        figura?
  \item dos \omterm{def:g6:lines:objects}{segmentos} de $8$ cm y $5$ cm, no \omterm{def:g6:lines:perp}{perpendiculares},
        que se cortan en su punto medio común. La figura que sale no es
        ninguna de las tres de la
        \cref{def:g6:shapes:quadrilaterals} --- describe lo que
        ves (¿qué lados parecen iguales? ¿paralelos?). Este
        \omterm{def:g4:geometry:polygon}{cuadrilátero} recibe su nombre en el \cref{ch:g7:central};
  \item resume las preguntas 1--4 en un diccionario de cuatro
        entradas: diagonales \emph{iguales / \omterm{def:g6:lines:perp}{perpendiculares} / las dos cosas /
        ninguna} (cortándose siempre en sus puntos medios)
        $\rightarrow$ nombre (o descripción) de la figura.
\end{enumerate}

\medskip
\noindent\textbf{Parte II --- Por qué funcionan las construcciones.}
\begin{enumerate}[resume]
  \item en la pregunta 2, las dos diagonales cortan la figura en cuatro
        \omterm{def:g6:shapes:triangles}{triángulos rectángulos} alrededor del punto de corte. ¿Cuáles son las
        longitudes de los dos catetos de cada uno de esos triángulos? Explica por qué
        los cuatro triángulos son copias exactas unos de otros y
        deduce por qué los cuatro lados de la figura son iguales --- el
        secreto del \cref{exo:g6:shapes:8};
  \item deduce una comprobación solo con compás: después de construir una figura
        cualquiera a partir de sus diagonales, ¿cómo puedes comprobar solo con el compás
        (sin medir con la regla) que los cuatro lados son
        iguales?
  \item escribe un programa de construcción
        (\cref{met:g6:shapes:program}) para un cuadrado cuyas
        \emph{diagonales} midan $6$ cm --- no sus lados;
  \item en el cuadrado de la pregunta 8, ¿es el lado más largo o más
        corto que $3$ cm, la \omterm{def:g3:division:half}{mitad} de la diagonal? Mide para
        averiguarlo. Ya se puede \emph{demostrar} la \omterm{def:g3:division:half}{mitad} de la
        respuesta: ir de una esquina a la siguiente pasando por el
        centro es un rodeo de $3 + 3$ cm, y un camino recto es
        siempre más corto que un rodeo --- así que el lado es menor
        que $6$ cm. (Demostrar que además supera los $3$ cm, y
        calcularlo exactamente, es tarea del
        \cref{ch:g8:pythagoras}.)
  \item una entrada nueva para el diccionario: traza un \omterm{def:g6:lines:objects}{segmento} $[AC]$ de
        $8$ cm; sobre él, marca el punto $H$ a $3$ cm de $A$;
        traza por $H$ el \omterm{def:g6:lines:objects}{segmento} perpendicular $[BD]$ de
        $5$ cm con $H$ como punto medio \emph{suyo}. Une $ABCD$. Esta
        figura es una \emph{cometa}\index{cometa} (una diagonal corta a
        la otra por la \omterm{def:g3:division:half}{mitad} y en \omterm{def:g3:shapes:rightangle}{ángulo recto}). Usando los triángulos de las esquinas como
        en la pregunta 6, halla qué lados de una \omterm{pb:g6:shapes:1}{cometa} son iguales a
        cuáles.
\end{enumerate}

\medskip
\noindent\textbf{Parte III --- Contar diagonales.}
Una \emph{diagonal} de un \omterm{def:g4:geometry:polygon}{polígono} une dos \omterm{def:g5:solids:def}{vértices} que no son
vecinos.
\begin{enumerate}[resume]
  \item dibuja un \omterm{def:g4:geometry:polygon}{cuadrilátero}, un \omterm{def:g4:geometry:polygon}{pentágono} (cinco \omterm{def:g5:solids:def}{vértices}) y un
        \omterm{def:g4:geometry:polygon}{hexágono} (seis \omterm{def:g5:solids:def}{vértices}), y traza todas sus diagonales.
        Cuéntalas: ¿cuántas hay en cada caso?
  \item contando sin dibujar, para el \omterm{def:g4:geometry:polygon}{hexágono}: desde un
        \omterm{def:g5:solids:def}{vértice}, ¿cuántas diagonales salen? (¿Con cuántos de los
        seis \omterm{def:g5:solids:def}{vértices} \emph{no} se puede unir mediante una diagonal?)
        Multiplicar por los seis \omterm{def:g5:solids:def}{vértices} cuenta cada diagonal
        exactamente dos veces --- ¿por qué? Recupera tu recuento de la pregunta
        11. (El mismo recuento doble apareció en los apretones de manos
        del \cref{pb:g6:wholes:1}.)
  \item usa el método de la pregunta 12 para contar las diagonales de
        un \omterm{def:g4:geometry:polygon}{polígono} de $10$ \omterm{def:g5:solids:def}{vértices} y después de uno de $20$ \omterm{def:g5:solids:def}{vértices};
  \item un \omterm{def:g4:geometry:polygon}{polígono} tiene exactamente $44$ diagonales. ¿Cuántos \omterm{def:g5:solids:def}{vértices}
        tiene? (Prueba el método de la pregunta 13 con unos cuantos
        candidatos.)
  \item comprueba tu método con el \omterm{def:g4:geometry:polygon}{polígono} más pequeño de todos: ¿cuántas
        diagonales tiene un triángulo y qué responde el
        método de recuento de la pregunta 12? Termina con una
        frase sobre lo que tienen en común los apretones de manos y las diagonales.
\end{enumerate}
\end{problem}
